\section*{Capítulo \ref{ch:g12:quantum-world} --- El mundo cuántico: fotones y niveles de energía}

\begin{solution}{exo:g12:quantum-world:1}
$E = hc/\lambda = \num{1.99e-25}/\num{5.30e-7} = \qty{3.75e-19}{J} =
\qty{2.35}{eV}$.
$N = \num{1.0e-3}/\num{3.75e-19} \approx \num{2.7e15}$ \omterm{def:g12:quantum-world:photon}{fotones} por
segundo.
\end{solution}

\begin{solution}{exo:g12:quantum-world:2}
$E = h\nu = 6.63\times10^{-34} \times 10^{8} = \qty{6.63e-26}{J}
\approx \qty{4.1e-7}{eV}$ --- unos cinco millones de veces menos que un
\omterm{def:g12:quantum-world:photon}{fotón} visible. Cualquier \omterm{def:g10:signals-and-waves:signal}{señal} de radio detectable involucra tantísimos
\omterm{def:g12:quantum-world:photon}{fotones} por segundo que el flujo resulta perfectamente liso.
\end{solution}

\begin{solution}{exo:g12:quantum-world:3}
$\nu_0 = W_0/h = \num{3.04e-19}/\num{6.63e-34} \approx
\qty{4.6e14}{Hz}$; $\lambda_0 = c/\nu_0 \approx \qty{654}{nm}$. A
\qty{633}{nm}, $E = 1240/633 = \qty{1.96}{eV} > \qty{1.9}{eV}$: sí, por
poco, con $E_{k,\max} \approx \qty{0.06}{eV} \approx \qty{1e-20}{J}$.
\end{solution}

\begin{solution}{exo:g12:quantum-world:4}
$E_3 - E_2 = -1.51 + 3.40 = \qty{1.89}{eV}$;
$\lambda = 1240/1.89 \approx \qty{656}{nm}$: la línea roja.
\end{solution}

\begin{solution}{exo:g12:quantum-world:5}
$\nu_0 \approx \qty{5.6e14}{Hz}$;
$W_0 = h\nu_0 \approx \qty{3.7e-19}{J} \approx \qty{2.3}{eV}$ (la
ordenada $-W_0$). La pendiente es la \omterm{def:g12:quantum-world:photon}{constante de Planck} $h$.
\end{solution}

\begin{solution}{exo:g12:quantum-world:6}
$E = 1240/400 = \qty{3.10}{eV}$;
$E_{k,\max} = 3.10 - 2.3 = \qty{0.80}{eV} = \qty{1.3e-19}{J}$;
$U_s = E_{k,\max}/e = \qty{0.80}{V}$.
\end{solution}

\begin{solution}{exo:g12:quantum-world:7}
Al duplicar la intensidad: la corriente se duplica (el doble de
\omterm{def:g12:quantum-world:photon}{fotones}) y $U_s$ no cambia (la misma \omterm{def:g10:energy-conservation:energy}{energía} por \omterm{def:g12:quantum-world:photon}{fotón}). Al subir la
\omterm{def:g12:oscillators-and-time:oscillator}{frecuencia}: $U_s$ sube, pues $E_{k,\max} = h\nu - W_0$. La imagen
ondulatoria ata la \omterm{def:g10:energy-conservation:energy}{energía} de los electrones a la intensidad y no a la
\omterm{def:g12:oscillators-and-time:oscillator}{frecuencia} --- exactamente lo que \emph{no} se observa.
\end{solution}

\begin{solution}{exo:g12:quantum-world:8}
Desde $n = 1$: \qty{13.6}{eV}, $\lambda = 1240/13.6 \approx
\qty{91}{nm}$, \omterm{def:g10:light-spectra:wavelength}{ultravioleta} lejano. Desde $n = 2$: \qty{3.40}{eV},
$\lambda \approx \qty{365}{nm}$, \omterm{def:g10:light-spectra:wavelength}{ultravioleta} cercano.
\end{solution}

\begin{solution}{exo:g12:quantum-world:9}
Desde el \omterm{def:g12:quantum-world:levels}{estado fundamental}, el salto más pequeño cuesta
$E_2 - E_1 = \qty{10.2}{eV}$ (\qty{122}{nm}): todas las líneas de
absorción son \omterm{def:g10:light-spectra:wavelength}{ultravioletas} y ninguna es visible. En la atmósfera
caliente de una estrella, los choques mantienen algunos \omterm{def:g11:fundamental-interactions:ladder}{átomos} en
$n = 2$, desde donde sí pueden absorberse los \omterm{def:g12:quantum-world:photon}{fotones} visibles de Balmer
(\qtyrange{1.89}{3.40}{eV}).
\end{solution}

\begin{solution}{exo:g12:quantum-world:10}
(a) $E_k = \qty{1.6e-17}{J}$,
$p = \sqrt{2m_e E_k} = \qty{5.4e-24}{kg.m/s}$,
$\lambda = h/p \approx \qty{1.2e-10}{m}$ --- del tamaño de un \omterm{def:g11:fundamental-interactions:ladder}{átomo}: los
cristales la difractan. (b) $p = \num{1.0e-4} \times 1.0 =
\qty{1.0e-4}{kg.m/s}$, $\lambda \approx \qty{6.6e-30}{m}$ --- $10^{20}$
veces menor que un \omterm{def:g11:fundamental-interactions:ladder}{átomo}: la \omterm{def:g10:signals-and-waves:wave}{onda} de la mosca es invisible para siempre.
\end{solution}

\begin{solution}{exo:g12:quantum-world:11}
Un \omterm{def:g12:quantum-world:photon}{fotón}: $\num{1.99e-25}/\num{5.05e-7} = \qty{3.94e-19}{J}$; noventa:
$\approx \qty{3.5e-17}{J}$. El mosquito:
$\tfrac12 \times \num{2.5e-6} \times 0.50^2 = \qty{3.1e-7}{J}$ ---
$10^{10}$ veces más. El ojo trabaja a un centenar de granos del límite
absoluto: casi un contador de \omterm{def:g12:quantum-world:photon}{fotones}.
\end{solution}

\begin{solution}{exo:g12:quantum-world:12}
$\nu_1 = \qty{7.41e14}{Hz}$, $\nu_2 = \qty{5.49e14}{Hz}$;
$h = e\,\Delta U_s/\Delta\nu = \num{1.60e-19} \times 0.79 /
\num{1.91e14} \approx \qty{6.6e-34}{J.s}$.
$W_0 = h\nu_2 - eU_{s,2} = \num{3.63e-19} - \num{0.64e-19} =
\qty{3.0e-19}{J} \approx \qty{1.9}{eV}$: cesio.
\end{solution}

\begin{solution}{exo:g12:quantum-world:13}
$1240/486 = \qty{2.55}{eV} = E_4 - E_2 = 3.40 - 0.85$: transición
$4 \to 2$. $1240/434 = \qty{2.86}{eV} = E_5 - E_2 = 3.40 - 0.54$:
transición $5 \to 2$.
\end{solution}

\begin{solution}{exo:g12:quantum-world:14}
Sección \omterm{def:g12:mechanical-waves:transverse}{transversal} $\pi r^2 \approx \qty{3.1e-20}{m^2}$, \omterm{def:g11:work-of-force:power}{potencia}
recogida $\qty{3.1e-26}{W}$;
$t = \num{3.7e-19}/\num{3.1e-26} \approx \qty{1.2e7}{s}$ --- cuatro
meses, frente a menos de un nanosegundo observado. La \omterm{def:g10:energy-conservation:energy}{energía} no puede
estar repartida por el \omterm{def:g12:mechanical-waves:wavefront}{frente de onda}: viaja concentrada, en granos.
\end{solution}

\begin{solution}{exo:g12:quantum-world:15}
$E_k = \qty{8.0e-16}{J}$,
$p = \sqrt{2m_e E_k} = \qty{3.8e-23}{kg.m/s}$,
$\lambda = h/p \approx \qty{1.7e-11}{m} = \qty{17}{pm}$ --- unas
$\num{3e4}$ veces más corta que \qty{550}{nm}. La \omterm{def:g12:light-as-wave:diffraction}{difracción} difumina
los detalles menores que $\sim\lambda$, así que encoger $\lambda$ en
$10^4$ compra $10^4$ en resolución.
\end{solution}

\begin{solution}{pb:g12:quantum-world:1}
\textbf{1.} $E = hc/\lambda = \num{1.99e-25}/\num{5.50e-7} =
\qty{3.62e-19}{J} = \qty{2.26}{eV}$.

\textbf{2.} $N = \num{1.0e3}/\num{3.62e-19} \approx \num{2.8e21}$
\omterm{def:g12:quantum-world:photon}{fotones} por segundo.

\textbf{3.} $P = 0.20 \times \qty{1.0e3}{W} = \qty{200}{W}$;
$E = 200 \times 5.0 = \qty{1.0}{kW.h}$ ($\qty{3.6e6}{J}$).

\textbf{4.} Con $10^{21}$ granos por segundo, el parpadeo estadístico es
fantásticamente pequeño: la granularidad se promedia y da una corriente
perfectamente lisa, igual que la arena vertida deprisa se derrama como
el agua.

\textbf{5.} El daño necesita un grano de \qtyrange{3}{4}{eV}: los granos
\omterm{def:g10:light-spectra:wavelength}{ultravioletas} cumplen, los \omterm{def:g10:light-spectra:wavelength}{infrarrojos} ($\sim \qty{0.1}{eV}$) no lo
hacen jamás --- el número no sustituye a la \omterm{def:g10:inertia:force}{fuerza} de cada grano.

\textbf{6.} $U_s$ es la tensión inversa que detiene justo a los
electrones más rápidos: $eU_s = E_{k,\max} = h\nu - W_0$.

\textbf{7.} $\nu = c/\lambda$: $8.22$, $7.41$, $6.88$ y $5.49$
$\times\,10^{14}\,\unit{Hz}$.

\textbf{8.} $eU_s = h\nu - W_0$ es afín en $\nu$ con pendiente $h$.
Pendientes sucesivas: $(0.94 - 0.37)/1.39 = 0.41$ y
$(1.50 - 0.94)/1.34 = 0.42$ \unit{eV} por $10^{14}\,\unit{Hz}$ --- los
tres puntos se alinean.

\textbf{9.} $h = e(U_{s,1}-U_{s,4})/(\nu_1-\nu_4) = \num{1.60e-19}
\times 1.13/\num{2.72e14} = \qty{6.64e-34}{J.s}$.

\textbf{10.} A \qty{546}{nm}:
$W_0 = h\nu - eU_s = \num{3.64e-19} - \num{0.59e-19} =
\qty{3.05e-19}{J} = \qty{1.90}{eV}$;
$\nu_0 = W_0/h = \qty{4.60e14}{Hz}$; $\lambda_0 \approx \qty{653}{nm}$.

\textbf{11.} Solo arranca electrones lo que tiene
$\lambda < \qty{653}{nm}$: el visible del violeta al rojo y el
\omterm{def:g10:light-spectra:wavelength}{ultravioleta}. Todo el lado \omterm{def:g10:light-spectra:wavelength}{infrarrojo} --- aproximadamente la mitad de la
\omterm{def:g10:energy-conservation:energy}{energía} del Sol --- le resulta inútil a este cátodo.

\textbf{12.} $(6.64 - 6.63)/6.63 \approx 0.2\,\%$ --- una medida de dos
por mil de una constante universal, hecha en un laboratorio escolar.

\textbf{13.} $E_2 = \qty{-3.40}{eV}$, $E_3 = \qty{-1.51}{eV}$,
$E_4 = \qty{-0.85}{eV}$, $E_5 = \qty{-0.54}{eV}$.

\textbf{14.} $3 \to 2$: \qty{1.89}{eV}; $4 \to 2$: \qty{2.55}{eV};
$5 \to 2$: \qty{2.86}{eV}.

\textbf{15.} $\lambda = 1240/E$: \qty{656}{nm} (rojo), \qty{486}{nm}
(azul verdoso), \qty{434}{nm} (violeta).

\textbf{16.} Las \omterm{def:g10:energy-conservation:energy}{energías} del \omterm{def:g11:fundamental-interactions:ladder}{átomo} están \omterm{def:g12:quantum-world:levels}{cuantizadas}, así que solo
puede emitir las diferencias discretas $E_p - E_n$: una escalera produce
líneas, no un arcoíris.

\textbf{17.} No: \qty{500}{nm} son \qty{2.48}{eV}, que no coinciden con
ninguna diferencia de niveles (desde $E_2$ hacia arriba salen $1.89$,
$2.55$, $2.86$ y \qty{3.40}{eV}; desde $E_1$ todo cuesta al menos
\qty{10.2}{eV}) --- el gas es transparente ahí.

\textbf{18.} $eU = \tfrac12 m_e v^2$:
$v = \sqrt{2 \times \num{1.60e-19} \times \num{1.0e4}/\num{9.11e-31}}
\approx \qty{5.9e7}{m/s} \approx 0.20\,c$ --- el tratamiento clásico es
aceptable.

\textbf{19.} $p = m_e v = \qty{5.4e-23}{kg.m/s}$;
$\lambda = h/p = \qty{1.2e-11}{m} = \qty{12}{pm}$.

\textbf{20.} Ganancia $= \num{5.50e-7}/\num{1.23e-11} \approx
\num{4.5e4}$. Las dos cifras del fin de semana: $h = \qty{6.64e-34}{J.s}$,
medida con un $0.2\,\%$ de error con unos \omterm{def:g11:color-light-sources:subtractive}{filtros} y un voltímetro, y un
microscopio cuarenta y cinco mil veces más nítido porque los electrones
ondulan.
\end{solution}
